% 第15章 量子反铁磁体：连续理论结果
\chapter{量子反铁磁体：连续理论结果}

第 12 章我们用 Haldane 映射把 $d$ 维量子 Heisenberg 反铁磁体（QHA）的长波关联与带附加 Berry 相位的 $d+1$ 维非线性 {$\sigma$} 模型（NLSM）的关联联系起来。第 13、14 章用连续近似求出了 $d=1,2,3$ 时 NLSM 的关联长度 $\xi$。

原则上，Haldane 映射在连续近似成立时有效，即

\begin{equation*}
\xi_{\mathrm{NLSM}} \gg \Lambda^ {- 1} \, , \quad \xi_{\mathrm{QHA}} \gg a ,
\tag{15.1}
\end{equation*}

$\Lambda$ 与 $a$ 分别是高动量截断与晶格常量。

在能用第 13、14 章的 NLSM 结果之前，必须理解 Berry 相位项 $e^{i\Upsilon}$ 对 (12.37) 右端的影响。

\section{一维：$\Theta$ 项}

考察一维链的 Berry 相位 $\Upsilon$（(12.28) 定义）。对缓变 $\hat{\mathbf{n}}(x)$，展开该项：

\begin{align*}
\Upsilon^ {d = 1} & = - S \sum_ {i} \left\{ \omega [ \hat {\mathbf {n}} ( x _ {2i} ) ] - \omega [ \hat {\mathbf {n}} ( x _ {2i - 1} ) ] \right\}\\
& = \frac {S}{2} \int \frac {d x}{a} \, \frac {\delta \omega}{\delta \hat {\mathbf {n}}} \cdot \partial_ {x} \hat {\mathbf {n}} \; a\\
& = 2 \pi S \, \Theta [ \hat {\mathbf {n}} ( x, \tau ) ] .
\tag{15.2}
\end{align*}

用 (10.34) 导出 $\Theta$ 的泛函形式：

\begin{equation*}
\Theta = \frac {1}{4 \pi} \int d \tau \int d x \, \left( \hat {\mathbf {n}} \times \partial_ {\tau} \hat {\mathbf {n}} \cdot \partial_ {x} \hat {\mathbf {n}} \right).
\tag{15.3}
\end{equation*}

$\Theta$ 是映射

\begin{equation*}
\hat {\mathbf {n}} : \left\{ \left[ - \frac {L}{2}, \frac {L}{2} \right), [ 0, \beta ) \right\} \rightarrow S ^ {2}
\tag{15.4}
\end{equation*}

的拓扑“绕数”或“Pontryagin 指数”，$S^{2}$ 是单位球面。$\Theta[\hat{n}]$ 是 $\hat{n}$ 值域的面积除以 $4\pi$。若定义域有周期边条件，$\hat{n}$ 必须把球面覆盖整数次，因此 $\Theta$ 必为整数。这个整数是拓扑稳定的：$\hat{n}$ 的任何连续形变都不能改变它。把 $\Theta$ 从 1 变到 0好比把包好的球面展开——除非在包装纸上开洞或撕裂，否则不可能。

$\Theta$ 项也出现在 Green 函数的路径积分表示中（那里 $\tau \to it$，见习题）。作为拓扑不变量，其变分为零，不影响经典运动方程。相位因子 $e^{i2\pi S\Theta}$ 对所有整数自旋为 1，但对半奇整数自旋 $S = \frac{1}{2}, \frac{3}{2}, \ldots$ 可正可负。这在路径积分中引起干涉效应，剧烈地改变基态关联与激发，使之不同于纯 NLSM。

存在交替 Berry 相位时连续理论难以求解。然而由 Lieb--Schultz--Mattis 定理（见 5.2 节）可知，半奇整数自旋链的最低激发能随格点数的倒数消失。由可积 $S = \frac{1}{2}$ 链的 Bethe 解，谱是无能隙的，基态关联的渐近衰减为

\begin{equation*}
\langle \mathbf {S} _ {i} \cdot \mathbf {S} _ {j} \rangle \sim ( - 1 ) ^ {( i - j )} \frac {\left( \ln | i - j | \right) ^ {\frac {1}{2}}}{| i - j |}.
\tag{15.5}
\end{equation*}

于是可以得出结论：$\Theta$ 项摧毁 Haldane 能隙并改变渐近关联\footnote{见 Affleck，及 Shankar 与 Read，文献注释。}。

现在给出拓扑干涉项与基态简并之间关系的一个启发性论证。考虑无自旋粒子在圆环上运动，环被磁通为 $\phi$ 的磁通管穿过，如图~\ref{fig:15.1}。哈密顿量简单地为

\begin{equation*}
\mathcal {H} = \frac {\hbar^ {2}}{2 m} \left( - i \partial_ {x} + \frac {\phi}{L \phi_ {0}} \right) ^ {2} ,
\tag{15.6}
\end{equation*}

$\phi_{0}=hc/me$ 是磁通量子。谱以整数 $n$ 标记（见图~\ref{fig:15.1}）：

\begin{equation*}
E ( n ) = \frac {\hbar^ {2}}{2 m L ^ {2}} \left( n + \frac {\phi}{\phi_ {0}} \right) ^ {2} \, , \quad \pm n = 0, 1, 2 \dots .
\tag{15.7}
\end{equation*}

\begin{figure}[H]
\centering
\includegraphics[width=0.5\textwidth]{fig15_1.jpg}
\caption{$\Theta$ 项的类比：环上粒子被磁通管穿过。半奇整数磁通量子导致基态简并。}
\label{fig:15.1}
\end{figure}

对整数 $\phi$，基态

\begin{equation*}
\left| n _ {0} \right\rangle = \left| - \phi / \phi_ {0} \right\rangle
\tag{15.8}
\end{equation*}

非简并。反之，对半奇整数 $\phi/\phi_{0}$，基态是两个简并态

\begin{equation*}
\left| - \phi / \phi_ {0} - \frac {1}{2} \right\rangle \, , \; \left| - \phi / \phi_ {0} + \frac {1}{2} \right\rangle .
\tag{15.9}
\end{equation*}

环上粒子与带 $\Theta$ 项的二维 NLSM 之间的类比，可从 (15.6) 的路径积分表示看出：

\begin{equation*}
Z = \int \mathcal {D} x \, e ^ {i \Upsilon [ x ( \tau ) ]} \exp \left( - \frac {m}{\hbar} \int_ {0} ^ {\beta} d \tau \, \frac {\dot {x} ^ {2}}{2} \right) ,
\tag{15.10}
\end{equation*}

$\Upsilon$ 是粒子与磁通管之间的纯规范耦合（即 Aharonov--Bohm 相位）：

\begin{align*}
\Upsilon & = \frac {2 \pi \phi}{L \phi_ {0}} \int_ {0} ^ {\beta} d \tau \, \dot {x}\\
& = 2 \pi n \frac {\phi}{\phi_ {0}} \, , \quad \pm n = 0, 1, 2 \dots .
\tag{15.11}
\end{align*}

$n$ 是 $x(\tau)$ 的绕数，即粒子绕环的圈数——它是 (15.3) 二维 Pontryagin 指数 $\Theta$ 的一维类比。于是对半奇整数磁通量子 $\phi$ 与半奇整数自旋 $S$，两个体系都在偶、奇拓扑扇区之间发生相消干涉。两个体系中，正是这种干涉造成基态简并。

\section{一维：整数自旋}

由 (15.2)，对整数自旋的一维量子 Heisenberg 反铁磁体可以忽略 Berry 相位因子 $e^{i\Upsilon}$。由 Haldane 映射，QHA 的连续近似是 $d=2$ NLSM，按 13.1 节它是经典 Heisenberg 模型的连续极限。由 Mermin--Wagner 定理 6.2，后者在一切有限温度下无长程序，翻译过来即：QHA 基态对一切 $S < \infty$ 无长程序。

关联长度 $\xi_{\mathrm{NLSM}}(f)$ 已在 (13.55) 求得；大 $N$ 近似 (14.26) 也给出同样结果。于是

\begin{align*}
\xi_{\mathrm{QHA}} ( d = 1, S ) & = \xi_{\mathrm{NLSM}} ( d = 2, f )\\
& \propto a \exp \left( \frac {2 \pi}{f} \right) \, , \quad S = 0, 1, \ldots .
\tag{15.12}
\end{align*}

由表 12.1，近邻模型的耦合常数为

\begin{equation*}
f = 2 / S ,
\tag{15.13}
\end{equation*}

从而

\begin{align*}
\xi_{\mathrm{QHA}} ( S, T = 0 ) & \approx a \exp ( \pi S ) ,\\
\Delta & \approx \frac {c}{a} \exp ( - \pi S )
\tag{15.14}
\end{align*}

$\Delta$ 是 (12.39) 给出的 Haldane 能隙。(15.14) 表明关联长度与 Haldane 能隙对半经典参数是非微扰的——不能按 $1/S$ 幂次从自旋波展开得到。形成对所有激发都有能隙的无序（自旋液体）基态是纯量子现象，没有经典类比；就此而言，它与单个粒子在能量势垒下的隧穿相似\footnote{隧穿率按 $\exp\left[-\frac{1}{\hbar}(\cdot)\right]$ 标度，$\hbar$ 是半经典参数。}。

回顾一维 AKLT Heisenberg 模型（见第 8 章，(8.26)）的基态具有指数衰减的关联 (8.30)；而且在单模近似（见 9.3 节）内我们曾在 (9.21) 估计其能隙。于是，以价键为基态的 $S = 1$ AKLT 模型与其他映射到 NLSM 的整数自旋 Heisenberg 反铁磁体定性相似。由于 AKLT 模型有依赖 $S$ 的相互作用，其在大 $S$ 下的标度行为不同于标准 Heisenberg 模型。

\section{二维}

零温时，正方格子上的 QHA 映射到三维 NLSM。按 (13.46)，后者在

\begin{equation*}
f _ {c} \geq 2 \pi^ {2}
\tag{15.15}
\end{equation*}

处有无序相。对近邻（nn）模型，Neves 与 Perez 证明 $S \geq 1$ 的基态有序。此外，级数展开与数值模拟表明 $S = \frac{1}{2}$ 的基态有序。

遗憾的是，推广 Heisenberg 模型的 $f(S,\Lambda)$ 的精确确定依赖短波正规化方案的细节（即“非普适”）。回忆连续近似本身依赖不等式 (12.10) 与 (15.16)。在临界温度 $T_{c} = f_{c}JS^{2}$ 附近，连近邻关联都大幅退化，在这一区域使用 NLSM 场论是有疑问的。

对连续近似的第一批修正涉及 Néel 场的点奇性。预期这些“刺猬”（hedgehog）在无序相中被热激发。Haldane 证明\footnote{见第 12 章文献注释。}刺猬引入非平凡 Berry 相位 $\Upsilon(2S \bmod 4)$，对 $S=\frac{1}{2},1$ 可能导致基态简并，而对例如 $S=2$ 则否。Read 与 Sachdev 的大 $N$ 途径与 Haldane 的方案一致，并预言由刺猬 Berry 相位引起的自旋 Peierls 有序\footnote{见第 18 章文献注释。}。

由 NLSM 在 $d+1$ 维转动下的对称性，对 $d$ 维动量截断 $\Lambda$，路径积分中最短时间尺度为 $2\pi/(c\Lambda)$。于是在

\begin{equation*}
T \ll c \Lambda
\tag{15.16}
\end{equation*}

区域，可以用宽 $c\beta$ 的有限板上 $d=3$ NLSM 的重整化群方程计算关联长度。Chakravarty、Halperin 与 Nelson（CHN）做到了 $\beta$ 函数的双圈阶。CHN 发现弱耦合区 $f$ 流向 Néel 不动点 $f_{c}=0$，基态有长程序。有限温下他们算得关联长度

\begin{equation*}
\xi_{\mathrm{QHA}} ^ {d = 2} = 0. 9 \, ( c / T ) \exp \left( \frac {2 \pi \rho_ {s}}{T} \right).
\tag{15.17}
\end{equation*}

可以把 $\xi_{\mathrm{QHA}}^{d=2}$ 与 (13.5)、(13.55) 给出的二维 Heisenberg 模型经典结果比较。量子涨落的实质效应是以常数 $Z_{\rho}$ 重整化劲度常数：

\begin{equation*}
\rho_ {s} = Z _ {\rho} \rho_ {s} ^ {cl}.
\tag{15.18}
\end{equation*}

(15.17) 的关联长度称为重整化经典关联。$S = 1/2$ 的 $Z_{\rho}$ 由级数展开算得为 $0.183$\footnote{见 R. R. P. Singh, Phys. Rev. B \textbf{39}, 9760 (1989)。}。

NLSM 对 $S = \frac{1}{2}$ QHA 的许多预言已在准二维 $La_{2}CuO_{4}$ 体系中得到数值与实验的证实。二维中可能的量子无序基态（量子自旋液体）的实验与理论面貌仍未定论。

\begin{chapbib}
\item 带 $\Theta$ 项的 NLSM 对半奇整数 $S$ 无能隙的证明：
  R. Shankar and N. Read, Nucl. Phys. B \textbf{336}, 457 (1990)。
\item 量子自旋链综述：
  I. Affleck, J. Phys. Condensed Matter \textbf{1}, 3047 (1989)。
\item 正方格子上 Heisenberg 反铁磁体 $S \geq 1$ 基态长程序的证明：
  E. J. Neves and J. F. Perez, Phys. Lett. A \textbf{114}, 331 (1986)。
\item 二维有限温关联的重整化群计算：
  S. Chakravarty, B. I. Halperin, and D. Nelson, Phys. Rev. B \textbf{39}, 2344 (1989)。
\item 关于自旋 1/2 二维 Heisenberg 反铁磁体与层状铜氧化物的两篇长篇综述：
  S. Chakravarty, ``High Temperature Superconductivity'', edited by K. S. Bedell, D. Coffey, D. E. Meltzer, D. Pines, and J. R. Schrieffer (Addison-Wesley, 1990)；
  E. Manousakis, Rev. Mod. Phys. \textbf{63}, 1 (1991)。
\end{chapbib}
